\section{Evidence, reproducibility, and formalization}

\subsection{Pinned inputs}

The repository snapshots are:
\begin{itemize}
\item \code{lsernau/grh-lean}:\\
\code{72fead102f663d5c7a407964a13dc0cf9d977207}.
\item \code{VladimirReshetnikov/ProveIt}:\\
\code{48efe71a726dddc5a505534edaa4c0fd190416b9}.
\end{itemize}
The directory inventory retains the incoming archive names, original Git blob identifiers and sizes. The two previously existing CM archives are included unchanged under \code{prior/}, together with their fresh replay logs. A separate source manifest records the selected upstream files and SHA-256 hashes. The package manifest covers all deliverable files.

\subsection{What was actually checked}

\begin{center}
\small
\begin{tabular}{p{.34\linewidth}p{.57\linewidth}}
\toprule
Item & Verification in this report\\
\midrule
Public Lean theorem scope & Pinned statement and source inspection; earlier external build reports attributed\\
Full upstream Lean proof & No new kernel replay; no local build claimed\\
Older four-term CM report & Complete supplied reproduction script rerun successfully\\
Current five-term CM report & Complete reproduction script rerun, including two last-quintic and independent sextic implementations\\
New ray-character minimum & Exact norm, form, finite-group and residue certificates; independently checked calculations\\
New confluent instances & Exact class-polynomial modular arithmetic with two recurrences, rank and integer determinants\\
General new lemmas & Written proofs with internal review; not Lean-verified\\
Literature priority & Not established; known ingredients attributed and new synthesis distinguished\\
\bottomrule
\end{tabular}
\end{center}

No finite numerical experiment is used to infer an all-exponent statement. The finite rank tests imply those statements through Theorem~\ref{thm:confluent}. Likewise, the minimum conductor calculation uses the proved ray-group quotient and an exhaustive finite list below the proposed minimum, plus a lower bound for all multi-prime supports.

\subsection{Running the package}

From the extracted archive directory, run
\begin{verbatim}
python3 reproduce.py
\end{verbatim}
This uses only the Python standard library for the new finite certificates. The generated JSON and text records are under \code{evidence/}. To replay the previous CM packages as well, use
\begin{verbatim}
python3 reproduce.py --include-prior
\end{verbatim}
The older verifiers also need their documented dependencies, including a C++17 compiler for the exhaustive residue sieve. Their original archives and reproduction files are retained rather than silently rewritten.

Rebuild this report with
\begin{verbatim}
latexmk -pdf -interaction=nonstopmode -halt-on-error report.tex
\end{verbatim}
The distributed PDF was compiled from the included source and rendered for layout review.

\subsection{A dependency-ordered Lean plan}

The smallest useful formalization starts with the finite-group extension criterion, keeping the constrained subgroup $D$ explicit. Next come the finite-free-algebra divided-difference lemma, the factorial-valuation estimate, and the transfer from a modular rank certificate to rational independence. These can be stated independently of the upstream automorphic machinery.

The CM ray sequence requires units, conjugation, finite primary components and global reciprocity. Its treatment of $p\mid|\mu(E)|$ should be formalized at the same time as the unrestricted case. The exact $-23$ calculation can then be attached through a small certified class-group presentation.

For the sparse examples, the modular certificates are small, but a fully Lean-certified singular-modulus theorem additionally needs the class-polynomial identities. Reusing their Python interval certificates is not itself a Lean proof. The dependency map under \code{formalization/} makes this boundary explicit.

The covolume rearrangement can be formalized as a short ordered-field lemma once the existing theorem is imported. Proving its arithmetic inputs for a new target remains the much larger task.
